% Revision table Table09. Requires \usepackage{booktabs,adjustbox}.
\begin{table}[t]
\centering
\caption{Predicted-class SHAP and LIME perturbation results at k=5 for selected BanglaBERT seed-42 examples.}
\label{tab:xai}
\footnotesize
\begin{adjustbox}{max width=\linewidth}
\begin{tabular}{lcccccc}
\toprule
Task & Method & n & Comp. drop & Suff. gap & Comp.-random & Bootstrap 95\% CI \\
\midrule
Fake News & SHAP & 20 & 0.466 & 0.118 & 0.239 & [0.096, 0.396] \\
Fake News & LIME & 20 & 0.434 & 0.287 & 0.207 & [0.042, 0.385] \\
Hate Speech & SHAP & 20 & 0.570 & 0.125 & 0.215 & [0.077, 0.366] \\
Hate Speech & LIME & 20 & 0.611 & 0.108 & 0.256 & [0.112, 0.412] \\
\bottomrule
\end{tabular}
\end{adjustbox}
\par\vspace{2pt}
\begin{minipage}{\linewidth}
\scriptsize\emph{Note.} Comprehensiveness measures the change in predicted-class probability after masking method-selected top-k tokens; the sufficiency gap measures the drop when retaining only them. Positive comp-minus-random is a comparative perturbation probe, not proof of semantic causality. These are 20 selected examples per task from a later measured protocol and must NOT be conflated with original manuscript n=141/200.
\end{minipage}
\end{table}
